\scanpage{47}

By the way, the fact that the $\pi_1$ for a surface are algebraic was pointed out (Tate tells me) by Hodge in Algebraic correspondances between surfaces, Proc. London Math. \unsure{Soc.} Seeies 2, Vol XLIV, 1938, p.226. It is rather striking that this statement should not have struck the algebraic geometers more, and has fallen into oblivion for \inserted{nearly} thirty years !

2) In car.0, $A(X)$ is known for $\dim X\leq4$. But $A(X)^0$ is not known if $\dim X=5$; the first interesting case would be for a variety $X\times Y$, $X$ of dim 3 and $Y$ a hyperplane section, as this would prove $C(X)$, see above.

Thus the main problems arise already for 1-cycles on threefolds, and partially even in car.0 . Urged by Kleimann's question, I \unsure{wihl|will} look again at my old scribbles on that subject (when I pretend to reduce the ``strong'' form of Lefschetz to the ``week'' one). As for the suggestion I made on the phone, to try to get any $X$ as \unsure{biratmonally|birationally} equivalent to a non singular $X'$, which is a spcialisation of a non singular $X''$, itself birationally equivalent to a non singular hypersurface -- this cannot work as Serre pointed out, because such an $X$ would have to be simply connected ! Thus if one wants to reduce somehow to the case of hypersurfaces, one will have to work also with singular ones, and see how to reformulate for singular varieties the standard conjectures \ldots

Sincerely your's

\mspage{46}
